\setlength{\footskip}{8mm}

\chapter{Experimental Setup}
\label{ch:experimental-setup}

\section{Purpose of the Experiments}

The experiments were designed as a case study of inference-time ML control. They examine how a pretrained latent diffusion model is conditioned by text, source latents, and gradients from a pretrained optical-flow model. Geometric initialization and restoration are evaluated as supporting stages around this ML process. The study does not establish broad statistical superiority, but it does include a controlled fixed-seed apple ablation over five variants and three seeds.

The evaluation therefore follows a staged structure:

\begin{enumerate}
  \item reproduce Stable Diffusion inference with RAFT-based Motion Guidance;
  \item replace file-only supervision with primitive-generated target flow;
  \item test latent gradient weighting, clipping, and geometric initialization;
  \item inspect saved flow loss, color loss, noise norm, and guidance norm trajectories;
  \item add source-hole restoration and sharp compositing;
  \item document which outputs are generated by the pipeline and which presets only display shipped assets.
\end{enumerate}

\section{Example Set}

The experiments use six main examples from the Motion Guidance-style dataset and repository assets: apple, teapot, lion, tree, woman, and topiary. These examples were selected because they cover both rigid translation and more difficult deformation cases.

\begin{table}[t]
  \caption{Examples used in the experimental setup.}
  \centering
  \small
  \begin{tabular}{p{0.9in}|p{1.15in}|p{1.25in}|p{1.55in}}
    \hline
    Example & Input Directory & Motion Type & Role in Evaluation \\ \hline \hline
    Apple & \path{./data/apple} & Translation & Included qualitative case study \\ \hline
    Teapot & \path{./data/teapot} & Translation & Documented; output not included \\ \hline
    Lion & \path{./data/lion} & Translation & Documented; output not included \\ \hline
    Tree & \path{./data/tree} & Stretch & Documented; output not included \\ \hline
    Woman & \path{./data/woman} & Non-rigid flow & Presentation preset; not evaluated \\ \hline
    Topiary & \path{./data/topiary} & Non-rigid flow & Presentation preset; not evaluated \\ \hline
  \end{tabular}
  \label{tab:example-set}
\end{table}

Only the apple example is used for the final metric table because it has the complete local input and output artifacts needed for repeated evaluation. Teapot, lion, and tree are documented configurations, but their complete repeated outputs are not part of the final ablation. Woman and topiary can load shipped reference assets and are therefore excluded from evaluation.

\section{Experimental Configurations}

\subsection{Pretrained Models}

Table~\ref{tab:ml-models} summarizes the learned models used during inference. No component is fine-tuned in this thesis.

\begin{table}[t]
  \caption{Pretrained ML components used by the editing pipeline.}
  \centering
  \small
  \begin{tabular}{p{1.2in}|p{1.35in}|p{2.45in}}
    \hline
    Component & Model & Function During Inference \\ \hline \hline
    Image representation & AutoencoderKL & Encodes RGB images to four-channel latents and decodes clean-latent estimates. \\ \hline
    Text conditioning & CLIP ViT-L/14 & Produces frozen prompt embeddings for cross-attention. \\ \hline
    Generative prior & Stable Diffusion v1.4 U-Net & Predicts conditional and unconditional diffusion noise. \\ \hline
    Motion estimator & RAFT & Predicts dense source-to-output flow inside the differentiable guidance loss. \\ \hline
    Learned restoration & LaMa, optional & Inpaints the source-hole mask after object translation. \\ \hline
  \end{tabular}
  \label{tab:ml-models}
\end{table}

Table~\ref{tab:experiment-configurations} summarizes the main configurations used in the final experimental setup. The target-flow file is still listed for primitive examples because it is used to localize the object support, while the actual displacement is generated from primitive parameters.

\begin{table}[t]
  \caption{Final experimental configurations.}
  \centering
  \small
  \begin{tabular}{p{0.7in}|p{1.0in}|p{1.0in}|p{1.05in}|p{1.2in}}
    \hline
    Example & Mask & Flow Support & Mode & Key Parameters \\ \hline \hline
    Apple & right mask & right flow & Primitive translate & $d_x=150$, $d_y=0$, $w_{out}=0.15$ \\ \hline
    Teapot & down mask & down flow & Primitive translate & $d_x=0$, $d_y=150$, $w_{out}=0.15$ \\ \hline
    Lion & left mask & left flow & Primitive translate & $d_x=-100$, $d_y=0$, $w_{out}=0.15$ \\ \hline
    Tree & squeeze mask & squeeze flow & Primitive stretch & $s_x=0.8$, $s_y=1.1$ \\ \hline
    Woman & hair mask & hair flow & Presentation asset mode & Excluded from evaluation \\ \hline
    Topiary & object mask & object flow & Presentation asset mode & Excluded from evaluation \\ \hline
  \end{tabular}
  \label{tab:experiment-configurations}
\end{table}

The documented inference settings were:

\begin{itemize}
  \item DDIM steps: $120$;
  \item classifier-free guidance scale: $7.5$;
  \item motion guidance weight: $30$;
  \item gradient clipping value: $60$;
  \item recursive steps: $1$;
  \item RAFT iterations: $1$;
  \item precision: fp32.
\end{itemize}

These values record the final apple ablation configuration. Because the study does not include a controlled hyperparameter search, they should not be interpreted as optimal settings.

\subsection{Recorded ML Diagnostics}

For every denoising step, the modified sampler records total guidance energy, RAFT flow loss, color-preservation loss, the norm of the U-Net noise prediction, and the norm of the applied guidance update. Intermediate decoded predictions and RAFT flow visualizations can also be saved every five steps. These diagnostics are intended to answer ML-specific questions: whether predicted motion approaches the target, whether appearance preservation dominates the energy, and whether the external gradient remains numerically stable.

\section{Compared Method Variants}

The evaluation compares the method in stages rather than only showing the final output. This is important because the thesis contribution is cumulative and because final compositing can hide the contribution of diffusion. The final apple ablation uses the following five variants:

\begin{table}[t]
  \caption{Method variants evaluated in the thesis.}
  \centering
  \small
  \begin{tabular}{p{1.35in}|p{3.6in}}
    \hline
    Variant & Description \\ \hline \hline
    Warp-inpaint-only & Diffusion is skipped; primitive translation, source-hole restoration, and final compositing are applied deterministically. \\ \hline
    Original Motion Guidance & Original file-flow Motion Guidance baseline using the prepared target flow. \\ \hline
    Diffusion, no RAFT & Primitive hard-warp initialization and preservation are used, but the RAFT guidance weight is zero. \\ \hline
    Diffusion + RAFT & Primitive hard-warp initialization and preservation are used with RAFT motion guidance and no final composite. \\ \hline
    Full pipeline & Primitive initialization, RAFT guidance, preservation, restoration, and final compositing are all enabled. \\ \hline
  \end{tabular}
  \label{tab:method-variants}
\end{table}

This staged description supports the component analysis in Chapter~\ref{ch:ablation}.

\section{Qualitative Evaluation Criteria}

The available thesis artifacts support both qualitative inspection and a focused quantitative apple ablation. The following criteria are used to interpret the saved metrics.

\subsection{Motion Consistency}

Motion consistency measures whether the estimated motion between the source image and edited image follows the target flow. Let $\Phi(x_0,x)$ be the RAFT-estimated flow between the source and output. The motion error is:

\begin{equation}
\mathcal{E}_{motion} =
\frac{1}{|M|}
\sum_{p \in M}
\left\|
\Phi(x_0,x)(p)-F(p)
\right\|_1.
\label{eq:motion-error}
\end{equation}

Lower motion error indicates that the output better follows the requested target motion.

\subsection{Source Preservation}

Source preservation measures how much the output changes outside the edited region. For complement mask $(1-M)$:

\begin{equation}
\mathcal{E}_{preserve} =
\frac{1}{|\bar{M}|}
\sum_{p \notin M}
\left\|
x(p)-x_0(p)
\right\|_1.
\label{eq:source-preservation-error}
\end{equation}

Lower preservation error means that the unedited background remains closer to the source image.

\subsection{Object Sharpness}

For translated objects, the object should remain sharp after movement. Sharpness is measured using the variance of the Laplacian inside the shifted support $S'$:

\begin{equation}
\mathcal{S}_{sharp} =
\mathrm{Var}\left(\nabla^2 x(p)\right), \quad p \in S'.
\label{eq:sharpness}
\end{equation}

Higher values generally indicate stronger high-frequency detail and less blur.

\subsection{Source-Hole Plausibility}

The source-hole region $H$ is the area exposed after the object moves. The evaluation checks whether this region becomes plausible background rather than retaining a ghost of the original object. A simple source-hole artifact score can be computed as:

\begin{equation}
\mathcal{E}_{hole} =
\frac{1}{|H|}
\sum_{p \in H}
\left\|
x(p)-B(p)
\right\|_1,
\label{eq:hole-score}
\end{equation}

where $B$ is the restored or expected background estimate. In the implemented metric script, related source-hole and boundary scores are computed using local gradient and color differences around the exposed region.

\subsection{Boundary Consistency}

Boundary consistency evaluates whether the object boundary contains visible seams. Let $\partial S'$ denote a narrow band around the shifted support boundary. A boundary seam score can be defined using gradient differences:

\begin{equation}
\mathcal{E}_{boundary} =
\frac{1}{|\partial S'|}
\sum_{p \in \partial S'}
\left\|
\nabla x(p)-\nabla B(p)
\right\|_1.
\label{eq:boundary-score}
\end{equation}

Lower values indicate smoother transitions between the composited object and the restored background.

\section{Evidence Categories}

To avoid overclaiming, outputs are separated into two categories:

\begin{itemize}
  \item \textbf{Evaluated pipeline output:} the apple ablation, consisting of five variants and three seeds.
  \item \textbf{Documented configurations:} teapot, lion, and tree settings recorded in development notes without complete local result artifacts.
  \item \textbf{Shipped presentation assets:} woman and topiary images loaded by an optional application mode; these are not experimental results.
\end{itemize}

This separation prevents interface demonstrations and incomplete records from being counted as empirical evidence.

\section{ML Environment and Validation}

The ML implementation uses Python, PyTorch, Stable Diffusion v1.4, RAFT, and optional LaMa inpainting. PyTorch autograd is required because the guidance energy must remain differentiable from RAFT's output to the diffusion latent. The core ablation commands are generated with:

\begin{verbatim}
python experiments/run_core_ablation.py \
  --results-root results/core_ablation \
  --seed <0|1|2> \
  --ddim-steps 120 \
  --raft-iters 1
\end{verbatim}

Metrics and plots are regenerated with:

\begin{verbatim}
python experiments/compute_metrics.py \
  --results-root results/core_ablation \
  --csv results/core_ablation/metrics.csv \
  --json results/core_ablation/metrics.json

python experiments/plot_diagnostics.py \
  --results-root results/core_ablation \
  --output-dir results/core_ablation/diagnostic_plots
\end{verbatim}

Runtime validation checks the Stable Diffusion checkpoint, RAFT checkpoint, input image, mask, target flow, inversion latent, and cached latents. The API and frontend are optional launch and visualization utilities; they are not part of the ML method.

\FloatBarrier
